\section{Slab 分配器概述}
\label{sec:slab-overview}

First-Fit 每次分配 $O(n)$ 遍历且易碎片化。\term{Slab 分配器}（源自 Solaris / Linux）
针对这两个弱点：按固定大小预切 slot，相同大小 $O(1)$ 取。

\subsection{7 级 Cache}

\begin{center}
\begin{tabular}{c c c l}
\toprule
编号 & 大小 & 每页 Slot 数 & 适用场景 \\
\midrule
0 & 32\,B  & $\approx 126$ & 小结构体、链表节点 \\
1 & 64\,B  & $\approx 63$  & 字符串缓冲区 \\
2 & 128\,B & $\approx 31$  & 进程控制块 \\
3 & 256\,B & $\approx 15$  & 中型数据 \\
4 & 512\,B & $\approx 7$   & 大结构体 \\
5 & 1024\,B & 3            & 页级缓冲 \\
6 & 2048\,B & 1            & 最大 cache \\
\bottomrule
\end{tabular}
\end{center}

\texttt{kmalloc(50)} $\to$ cache-64，内部碎片 = $14$\,B。

\subsection{Slab 内存布局}

每个 slab 占一页（4\,KiB），页头放 \texttt{slab\_header\_t}（含 128 位 bitmap）：

\begin{figure}[H]
  \centering
  % figures/ch10/fig07-slab-layout.tex — auto-extracted from chapter source
\begin{tikzpicture}[cell/.style={draw, minimum height=0.7cm, font=\ttfamily\small, anchor=west}]
  \node[cell, minimum width=2.2cm, fill=red!12]   at (0, 0) {\small slab\_hdr};
  \node[cell, minimum width=1.4cm, fill=green!12]  at (2.2, 0) {\small obj[0]};
  \node[cell, minimum width=1.4cm, fill=green!12]  at (3.6, 0) {\small obj[1]};
  \node[cell, minimum width=1.4cm, fill=green!12]  at (5.0, 0) {\small obj[2]};
  \node[font=\small] at (7.1, 0) {$\cdots$};
  \node[cell, minimum width=1.4cm, fill=green!12]  at (7.6, 0) {\small obj[N]};
  \draw[{|}-{|}, thick] (0, -0.6) -- (2.2, -0.6);
  \node[font=\scriptsize, below] at (1.1, -0.6) {$\sim$40\,B};
  \draw[{|}-{|}, thick] (0, 0.9) -- (9.0, 0.9);
  \node[font=\scriptsize, above] at (4.5, 0.9) {一页 = 4096\,B};
\end{tikzpicture}

  \caption{Slab 内存布局}
  \label{fig:slab-layout}
\end{figure}

分配：\texttt{bitmap\_find\_first\_set} $\to$ 清位 $\to$ 计算地址，$O(1)$。
释放：\texttt{ptr \& \~{}0xFFF} 找到页头 header，算 slot index $\to$ 置位，$O(1)$。

\begin{designbox}
\texttt{slab\_header\_t} 放在页头而非单独分配的三个原因：
① 避免鸡生蛋——slab 自己就是分配器，metadata 不能再 \texttt{kmalloc}；
② \texttt{kfree} 时页对齐即得 header，$O(1)$；
③ 每页只损失 $\sim$40\,B（$<1\%$）。
\end{designbox}

% ============================================================================

